\section{\eps: Definition and Construction}
\label{sec:epsilon}

\subsection{Per-token entropy}
\label{subsec:per-token}

For each generated token $t$, the API returns top-$k$ candidates with
log-probabilities $\ell_1, \dots, \ell_k$. Let
$p_j = \exp(\ell_j) / \sum_j \exp(\ell_j)$ be the renormalized
distribution. Define
\begin{equation}
    \eps_t \;=\; -\frac{1}{\log k} \sum_{j=1}^{k} p_j \log p_j \;\in\; [0,1].
    \label{eq:eps}
\end{equation}
Normalization by $\log k$ maps a one-hot distribution to $0$ and a
uniform distribution to $1$. We use $k=5$ (the OpenAI cap), which
covers the dominant mass for code tokens; tails beyond top-5 are
negligible for the binary flag decision. The name is
deliberate: in numerical analysis, $\eps$ denotes a small perturbation
that propagates through a computation --- a high-\eps token is such
a perturbation at a specific decision point.

\subsection{Three types of code-generation uncertainty}
\label{subsec:types}

Raw per-token entropy, applied naively, produces too many false positives
to be useful: the bulk of high-entropy tokens in generated code are not
risk indicators. Understanding which ones are, and why, is what drives
the filtering design in the rest of this section. Empirical analysis of
our calibration benchmark reveals three categorically distinct types of
token-level uncertainty in generated code.

\paragraph{Type 1 --- Cosmetic (naming).} Uncertainty about what to
\emph{call} something (\texttt{sort\_list} vs.\ \texttt{sort\_numbers}).
$\eps$ reaches $0.7$--$0.9$, but any choice is correct. Dominates LOW
prompts: the top flag on $9$ of $10$ LOW prompts was a function or
parameter name on the first line.

\paragraph{Type 2 --- Consequential (API/library decisions).}
Uncertainty about which \emph{API or library} to use. Example:
\texttt{stripe.Charge.create()} (prob.\ $0.52$) vs.\
\texttt{stripe.PaymentIntent.create()} (prob.\ $0.47$), $\eps = 0.73$
at the attribute token. This is the target signal.

\paragraph{Type 3 --- Implementation-approach.} Uncertainty among
multiple correct algorithms (recursive vs.\ iterative Fibonacci,
for-loop vs.\ comprehension). The residual false-positive floor on
simple utility prompts.

Only Type 2 correlates with functional failure. Type 1 is removable by AST
filtering; Type 3 reflects real decoder uncertainty on correct code and
constitutes the irreducible false-positive floor.

\begin{table}[h]
\centering
\small
\begin{tabular}{c p{2.6cm} p{7cm} p{2.4cm} p{1.2cm}}
\toprule
\textbf{Type} & \textbf{Label} & \textbf{Source} & \textbf{\eps fires?} & \textbf{Risk} \\
\midrule
1 & Cosmetic & Naming choices (\texttt{sort\_list} vs \texttt{sort\_numbers}) & Yes (0.7--0.9) & None \\
\multicolumn{5}{l}{\quad\small\textit{Filter: AST declaration filter removes this uncertainty before scoring.}} \\[4pt]
2 & Consequential & API/library decisions (\texttt{Charge} vs \texttt{PaymentIntent}) & Yes (0.5--0.9) & \textbf{High} \\
\multicolumn{5}{l}{\quad\small\textit{Note: Target signal --- this is what the system is designed to catch.}} \\[4pt]
3 & Implementation & Algorithm structure (recursive vs.\ iterative) & Yes (0.3--0.7) & None \\
\multicolumn{5}{l}{\quad\small\textit{Note: Irreducible FP floor; rare in production API-integration code.}} \\
\bottomrule
\end{tabular}
\caption{The three-type uncertainty taxonomy. Type 2 is the target signal; Type 1 is removed by AST filtering; Type 3 is the irreducible false-positive floor.}
\label{tab:taxonomy}
\end{table}

With this taxonomy in hand, the filter design follows directly: remove
Type 1 syntactically, leave Type 3 as a known floor, and isolate Type 2.

\subsection{AST filtering}
\label{subsec:ast}

We eliminate Type 1 via the generated code's AST. After stripping
markdown fences and parsing with Python's \texttt{ast}, we identify
\emph{declaration positions} --- source ranges where the model is
choosing a name rather than an API.
\textbf{Filtered:} function/method names, parameter names
(\texttt{ast.arg}), local variable names (\texttt{ast.Name} in
\texttt{Store}).
\textbf{Retained:} import targets, attribute accesses, keywords, type
annotations. Token-to-declaration matching uses range overlap (BPE
introduces a leading space); a pure-identifier guard prevents over-%
filtering of delimiter-prefixed tokens such as \texttt{(db}.

\begin{table}[h]
\centering
\small
\begin{tabular}{lc}
\toprule
\textbf{Configuration} & \textbf{LOW FP rate} \\
\midrule
Naive token entropy (no filtering)                & $90\%$ \\
+ Whitespace / comment filtering                   & $90\%$ \\
+ System prompt naming conventions                 & $60\%$ \\
+ AST declaration filtering                        & $60\%$ \\
+ Fence-format stripping                           & $50\%$ \\
\midrule
Residual (Type 3 only)                             & $50\%$ \\
\bottomrule
\end{tabular}
\caption{Successive filtering stages on 10 LOW prompts. Naming-%
convention prompt and AST filter are partially overlapping; without
the prompt, AST filtering alone drops FP from $90\%$ to $70\%$.
Residual $50\%$ is Type 3, rare in production API-integration code.}
\label{tab:filtering}
\end{table}

\subsection{Aggregation: max vs.\ mean}
\label{subsec:aggregation}

We aggregate with
$\eps_\text{file} = \max\{\eps_t : t \text{ retained}, \eps_t > \delta\}$,
$\delta = 0.30$. The compound alternative
$1 - \prod(1 - \eps_t)$ converges to $1.0$ for any non-trivial function.
Across all $90$ runs of the benchmark ($3 \times 30$), max and mean
produced identical fire/no-fire decisions at every threshold; we prefer
max because it localizes uncertainty to a specific actionable token. The
contrast with the sequence-level mean token probability baseline
(\papg) is deliberate: averaging over every token dilutes a single
uncertain decision; the max after AST filtering does not. We revisit
this choice empirically in \S\ref{sec:comparison}.

\subsection{Thresholds and tiers}
\label{subsec:tiers}

Three tiers: \textsc{complete} ($\eps < 0.30$), \textsc{flagged}
($0.30 \le \eps < 0.95$), \textsc{paused} ($\eps \ge 0.95$). The load-%
bearing boundary is the binary \textsc{complete} vs.\ \textsc{flagged+}:
everything \textsc{flagged+} is forwarded to the stage-2 reviewer.
\textsc{paused} exists as a design boundary for genuinely extreme
uncertainty; whether it holds code delivery or simply runs the review
loop like \textsc{flagged} is an implementation choice discussed in
\S\ref{sec:discussion}. These are UX signals, not calibrated
probabilities --- the \textsc{flagged} empirical failure rate is
$\sim 27\%$ (\S\ref{sec:evaluation}); \textsc{paused} fires rarely
enough in our evaluation that its failure rate is a small-sample
estimate. In our qualitative scenarios A--D across four models, no
HIGH prompt crossed the \textsc{paused} threshold.

\subsection{Multi-model calibration at a glance}
\label{subsec:calibration-preview}

\eps at the \textsc{flagged} threshold fires on HIGH prompts (API/library-%
selection tasks where version splits exist and an incorrect choice
causes functional failure) much more often than on LOW prompts (pure-%
logic tasks --- sort, palindrome, Fibonacci --- where any correct
algorithm is functionally equivalent). We define these two categories
here because they are referenced throughout the evaluation; a detailed
calibration table appears in \S\ref{subsec:calibration}. Across four
models (GPT-4o, GPT-4o-mini, GPT-4-turbo, DeepSeek V3), \eps achieves
$70$--$85\%$ HIGH detection at $50$--$70\%$ LOW false-positive rate.

\subsection{Adaptive ensemble}
\label{subsec:ensemble}

The production implementation refines tier placement within
\textsc{flagged+} using a rolling mean-\eps baseline and a top-$1$
logprob check; neither affects the binary \textsc{flagged+} gate.
Conformal calibration is left to future work.
